% !PRINTED 653
\weberpage{653}

These two numbers admit, as follows immediately from (7), the substitutions $(x_0, x_\varkappa)$, and are therefore numbers of the field $\Omega_m$.

If $n$ is not divisible by $q$, then on account of the irreducibility of the cyclotomic equation none of the resolvents $(r^n, x_0)$ can vanish unless they all vanish; for the equation
\[
(r, x_0)^m = 0
\]
would be an equation for $r$ with rational coefficients which admits all the substitutions $(r, r^n)$.

In the same way none of the resolvents $(r^s, x_0)$ can vanish unless the whole family vanishes in which $s$ has a definite greatest common divisor with $m$.

If all the $(r, x_0)$ vanish, formula (9) is not applicable. But then we can consider directly the resolvents $(r^q, x_0)$, which belong to a field of degree $m q^{-1}$ arising from
\[
y_0 = x_0 + x_q + x_{2q} + \cdots + x_{m-q}
\]
and if these quantities also vanish, we pass on to the resolvents $(r^{q^2}, x_0)$, and so on.

But if $(r, x_0)$ does not vanish, then by formula (9) all the $(r^s, x_0)$ can be represented rationally in terms of $(r, x_0)$ and of cyclotomic numbers.

Accordingly it is sufficient for our purpose if we can prove that the \gesperrt{non-vanishing resolvents} $(r, x_0)$ \gesperrt{are cyclotomic numbers}.

If we let $n$ run through a complete system of incongruent numbers not divisible by $q$, then
\[
\sum n \equiv 0 \ (\mathrm{mod}\ m),
\]
for the numbers $n$ split into pairs of numbers completing one another to $m$. From this it follows by (7) that the product
\weq{10}{\prod^n (r^n, x_0) = a}
is a number in $\Omega_m$. But since this number, moreover, does not change when any one of the substitutions $(r, r^n)$ is carried out in it, $a$ is a \gesperrt{rational number}.
